% Part: first-order-logic
% Chapter: natural-deduction
% Section: soundness-identity

\documentclass[../../../include/open-logic-section]{subfiles}

\begin{document}

\olfileid{fol}{ntd}{sid}

\olsection{\usetoken{S}{identity}-সহ বিশুদ্ধতা}

\begin{prop}
$\eq$-এর বিধিসহ স্বাভাবিক নিষ্পাদন বিশুদ্ধ।
\end{prop}

\begin{proof}
$\eq[t][t]$ রূপের যেকোনো !!{formula} বৈধ, কারণ প্রত্যেক
!!{structure}~$\Struct M$-এর জন্য $\Sat{M}{\eq[t][t]}$। (মনে করো,
পদ~$t$-কে বদ্ধ ধরা হয়েছে, অর্থাৎ এতে কোনো চলরাশি নেই; তাই চলরাশি-আরোপ
এখানে অপ্রাসঙ্গিক।)

ধরা যাক কোনো !!a{derivation}-এর শেষ অনুমান হলো \Elim{\eq}; অর্থাৎ
!!{derivation}-টির রূপ নিচের মতো:
\begin{prooftree}
  \AxiomC{$\Gamma_1$}
  \RightLabel{$\delta_1$}
  \DeduceC{$\eq[t_1][t_2]$}
  \AxiomC{$\Gamma_2$}
  \RightLabel{$\delta_2$}
  \DeduceC{$!A(t_1)$}
  \RightLabel{\Elim{\eq}}
  \BinaryInfC{$!A(t_2)$}
\end{prooftree}
পূর্বধারণা $\eq[t_1][t_2]$ ও $!A(t_1)$ যথাক্রমে !!{undischarged}
অনুমিতি~$\Gamma_1$ ও $\Gamma_2$ থেকে !!{derive}d। আমরা দেখাতে চাই,
$\Gamma_1 \cup \Gamma_2$ থেকে~$!A(t_2)$ অনুসৃত হয়। এমন কোনো
!!a{structure}~$\Struct{M}$ বিবেচনা করো যেখানে $\Sat{M}{\Gamma_1 \cup
  \Gamma_2}$। আরোহের অনুমান থেকে $\Sat{M}{!A(t_1)}$ এবং
$\Sat{M}{\eq[t_1][t_2]}$। অতএব $\Value{t_1}{M} = \Value{t_2}{M}$।
$s$ যেকোনো চলরাশি-আরোপ এবং $m = \Value{t_1}{M} = \Value{t_2}{M}$
ধরি। \olref[fol][syn][ext]{prop:ext-formulas} অনুযায়ী,
$\Sat{M}{!A(t_1)}[s]$ যদি এবং কেবল যদি
$\Sat{M}{!A(x)}[\Subst{s}{m}{x}]$ যদি এবং কেবল যদি
$\Sat{M}{!A(t_2)}[s]$। যেহেতু $\Sat{M}{!A(t_1)}$, তাই
$\Sat{M}{!A(t_2)}$।
দ্বিতীয় অভিন্নতা-অপসারণ বিধির ক্ষেত্রটি পদদুটির ভূমিকা অদলবদল করে
একইভাবে প্রমাণিত হয়।
% BN-SRC-041 TARGET-CORRECTION-BEGIN
\par\smallskip\noindent\textnormal{\small\textbf{উৎস-সংশোধন (BN-SRC-041):} আগের অংশে অভিন্নতার দুটি অপসারণ-বিধি দেওয়া হলেও হিমায়িত ইংরেজি বিশুদ্ধতা-প্রমাণটি শুধু $!A(t_1)$ থেকে $!A(t_2)$ অভিমুখটি যাচাই করে শেষ হয়েছে। বাংলা প্রমাণে পদদুটির ভূমিকা অদলবদল করে দ্বিতীয় অভিমুখটির একই যুক্তি স্পষ্ট করা হয়েছে।} % BN-SRC-041 TARGET-CORRECTION-END
\end{proof}

\end{document}
